\lecture{10}{Tue 31 Mar 2026 11:03}{owo}

\begin{theorem}[Engel]\label{thm:hi}
	$\mathfrak{g} $ is nilponent if and only if $\mathrm{ad} : \mathfrak{g}  \to \End(\mathfrak{g} ) $ is strictly upper triangular. Thus any nilpotent representation is strictly upper triangular.
\end{theorem}
\begin{theorem}[Lie]\label{thm:Lie}
	$\mathfrak{g} $ is solvable if and only if $\mathrm{ad} : \mathfrak{g}  \to \End(\mathfrak{g} ) $ is upper triangular. Thus any solvable representation is upper triangular.
\end{theorem}

A Killing form $K(X,Y) = \mathrm{tr} (\mathrm{ad}_X, \mathrm{ad} _Y)$ on $\mathbb{C} $ is a symmetric bilinear form on $\mathfrak{g} $. It may be degenerate, meaning $K(X,-)=0$ even if $X\in Z(\mathfrak{g} )$. Cartan has a theorem relating semisimplicity of Lie algebras to their Killing forms.

\begin{theorem}[Cartan]\label{thm:waow}
	A Lie algebra $\mathfrak{g} $ is semi-simple if and only if its Killing form is nondegenerate. To prove this, we must also prove: $\mathfrak{g} $ is sovable if and only if $K(\mathfrak{g} ,[\mathfrak{g} ,\mathfrak{g} ])=0$.
\end{theorem}
